% ============================================================
%  USPSC-Cap3-Metodologia_3.2-Fonte_de_Dados_e_Coleta.tex
%  Seção 3.2 — Fonte de Dados e Processo de Coleta
% ============================================================

\section{Fonte de Dados e Processo de Coleta}
\label{sec:fonte_dados}

\subsection{Justificativa da Escolha do X/Twitter como Fonte}
\label{subsec:justificativa_twitter}

O X/Twitter foi selecionado como fonte primária de dados por congregar características estruturais que o tornam especialmente adequado para a análise de sentimento aplicada ao mercado financeiro. Em primeiro lugar, a plataforma opera como canal de comunicação em tempo real, com publicações cronológicas e públicas que registram o humor e as expectativas dos agentes financeiros de forma praticamente instantânea \cite{bollen2011}. Em segundo lugar, a arquitetura da plataforma — baseada em publicações curtas (\textit{tweets}) com limite de caracteres — produz textos concisos e com
orientação semântica, mais adequados à tarefa de classificação de polaridade do que textos longos de natureza mais ambígua \cite{ranco2015}.

A literatura internacional consolidou o X/Twitter como objeto de pesquisa em finanças, desde os trabalhos pioneiros de \citeonline{bollen2011} e \citeonline{antweiler2004}, que demonstraram a relação entre o conteúdo das publicações e o comportamento de mercado. No contexto brasileiro, trabalhos como os de \citeonline{sousa2025} e \citeonline{falcao2020} também recorreram à plataforma como fonte de dados para análise de sentimento financeiro.

Em comparação com alternativas como artigos completos de veículos de imprensa, as publicações do X/Twitter apresentam duas vantagens metodológicas relevantes: (i) maior facilidade de extração estruturada de dados via API, dispensando técnicas complexas de \textit{web scraping}; e (ii) formato textual conciso e objetivo, mais adequado para análise de polaridade semântica no nível de documento.

\subsection{Veículos Selecionados}
\label{subsec:veiculos}

Como \textit{proxy} para o sentimento do mercado financeiro brasileiro, foram selecionadas as publicações de dois veículos especializados em economia e finanças com presença ativa na plataforma X/Twitter: \textbf{InfoMoney} (\texttt{@InfoMoney}) e \textbf{Investing Brasil} (\texttt{@InvestingBrasil}).

A escolha desses veículos é fundamentada em três critérios:

\begin{enumerate}
	\item \textbf{Relevância editorial}: ambos são referências reconhecidas de jornalismo econômico-financeiro no Brasil, com cobertura sistemática de eventos macroeconômicos, decisões de política monetária, resultados corporativos e movimentações do mercado de capitais nacional;

	\item \textbf{Volume e regularidade de publicações}: as contas institucionais mantêm frequência elevada e contínua de publicações, garantindo densidade suficiente de dados para a construção de séries temporais de sentimento;

	\item \textbf{Consistência do público-alvo}: ao contrário de contas de pessoas físicas, cujo conteúdo pode ser heterogêneo, os veículos especializados produzem publicações com foco temático predominantemente financeiro, reduzindo a necessidade de filtragem por relevância.
\end{enumerate}

Reconhece-se que a opção por veículos institucionais em detrimento de contas de investidores individuais representa uma escolha metodológica deliberada, que delimita o escopo do sentimento capturado ao discurso editorial — e não ao sentimento difuso dos participantes de varejo do mercado. Essa limitação circunscreve o alcance das generalizações admissíveis e é reconhecida como restrição metodológica deste trabalho.

\subsection{Período de Análise}
\label{subsec:periodo}

O corpus compreende publicações realizadas entre \textbf{17 de outubro de 2025} e \textbf{12 de fevereiro de 2026}, totalizando aproximadamente quatro meses de dados — intervalo determinado pelo registro mais antigo e pelo mais recente efetivamente presentes na base coletada via API. Esse recorte temporal foi determinado por um conjunto de restrições metodológicas e operacionais.

Do ponto de vista \textbf{técnico}, o acesso à API do X/Twitter está condicionado a planos comerciais com janelas históricas limitadas. A recuperação retroativa de grandes volumes de dados requer subscrição ao plano \textit{Pro} ou \textit{Enterprise} da API, cujos custos são proibitivos para pesquisas acadêmicas sem financiamento dedicado. O período analisado corresponde à janela de dados coletados em tempo real durante a execução da pesquisa, garantindo acesso completo e sem truncamento ao histórico de publicações.

Do ponto de vista \textbf{analítico}, o período de aproximadamente quatro meses abrange um conjunto relevante de eventos macroeconômicos do cenário brasileiro — incluindo reuniões do Comitê de Política Monetária (COPOM), divulgação de indicadores de inflação (IPCA e INPC), dados de atividade econômica e movimentações relevantes nos preços dos ativos — o que possibilita a avaliação qualitativa da coerência entre os índices de sentimento extraídos e o contexto econômico-financeiro do período.

\subsection{Processo de Coleta via API}
\label{subsec:coleta_api}

\subsubsection{\textit{Endpoint} Utilizado e Autenticação}
\label{subsubsec:endpoint}

A coleta foi realizada por meio do \textit{endpoint} \textbf{GET \texttt{/2/users/\{id\}/tweets}} da API v2 do X/Twitter\footnote{Documentação oficial: \url{https://developer.twitter.com/en/docs/twitter-api/tweets/timelines/api-reference/get-users-id-tweets}. Acesso em: 12 fev. 2026.}, que retorna a linha do tempo de publicações de um usuário específico identificado pelo seu \texttt{id} numérico. A opção por esse \textit{endpoint} — em vez do \textit{endpoint} de busca (\textit{Search}) — é metodologicamente justificada pelo escopo da coleta: como os veículos de interesse (InfoMoney e Investing Brasil) são previamente conhecidos e possuem identificadores únicos fixos, a recuperação direta pela linha do tempo do usuário é mais eficiente e determinística do que uma busca por palavras-chave, eliminando a necessidade de filtros adicionais de autoria.

A autenticação foi realizada via \textbf{Bearer Token} (\textit{app-only authentication}), mecanismo que concede acesso de leitura ao conteúdo público da plataforma sem requerer autorização individual de cada conta coletada. O token foi armazenado como variável de ambiente (\texttt{TWITTER\_ACCESS\_TOKEN}) e transmitido no cabeçalho HTTP de cada requisição no formato \texttt{Authorization: Bearer <token>}.

\subsubsection{Parâmetros da Requisição}
\label{subsubsec:parametros_api}

Para cada veículo monitorado, as requisições foram parametrizadas conforme o \autoref{qua:params_api}. O intervalo temporal foi delimitado pelos campos \texttt{start\_time} e \texttt{end\_time}, ambos no formato ISO~8601 (\texttt{YYYY-MM-DDTHH:mm:ssZ}, UTC). O volume máximo de resultados por página foi fixado em 100 registros — limite superior permitido pelo \textit{endpoint} —, e a paginação foi controlada pelo parâmetro \texttt{pagination\_token}, utilizado de forma iterativa até o esgotamento dos resultados disponíveis.

\begin{quadro}[htb]
	\caption{\label{qua:params_api}Parâmetros utilizados nas requisições ao \textit{endpoint} \texttt{GET /2/users/\{id\}/tweets}}
	
	\begin{tabular}{|p{4.5cm}|p{10.5cm}|}
		\hline
		\textbf{Parâmetro} & \textbf{Valor / Descrição} \\
		\hline
		
		\texttt{id} (path) & Identificador numérico do veículo: \texttt{@InfoMoney} e \texttt{@InvestingBrasil}, passados dinamicamente a cada ciclo de coleta \\
		\hline
		
		\texttt{start\_time} & Limite inferior do intervalo temporal (ISO 8601, UTC) \\
		\hline
		
		\texttt{end\_time} & Limite superior do intervalo temporal (ISO 8601, UTC) \\
		\hline
		
		\texttt{max\_results} & \texttt{100} (máximo permitido pelo \textit{endpoint}) \\
		\hline
		
		\texttt{tweet.fields} & \texttt{created\_at}, \texttt{note\_tweet}, \texttt{author\_id}, \texttt{public\_metrics}, \texttt{lang}, \texttt{source}, \texttt{entities}, \texttt{context\_annotations}, \texttt{geo} \\
		\hline
		
		\texttt{user.fields} & \texttt{name}, \texttt{username}, \texttt{location}, \texttt{description}, \texttt{public\_metrics} \\
		\hline
		
		\texttt{pagination\_token} & Token retornado pela API para recuperação da próxima página de resultados; omitido na primeira requisição de cada ciclo \\
		\hline
	\end{tabular}
	\legend{Fonte: Elaborado pelo autor com base na documentação da API v2 do X/Twitter e no código-fonte do \textit{pipeline} de coleta.}
\end{quadro}

Um aspecto técnico relevante diz respeito ao campo \texttt{note\_tweet}: publicações de usuários com assinatura X Premium podem ultrapassar o limite convencional de 280 caracteres por meio desse campo, que armazena o texto estendido separadamente do campo \texttt{text} padrão. O \textit{pipeline} implementado prioriza o conteúdo de \texttt{note\_tweet} quando disponível, recorrendo ao \texttt{text} convencional apenas na sua ausência — decisão que garante a recuperação do conteúdo textual integral das publicações. Registros sem conteúdo textual em nenhum dos dois campos são descartados durante a coleta.

\subsubsection{\textit{Pipeline} de Coleta e Armazenamento}
\label{subsubsec:pipeline_coleta}

O \textit{pipeline} de coleta é a primeira etapa (\texttt{pipeline/01\_extraction/}) de uma estrutura modular de seis estágios numerados. Os módulos que compõem essa etapa, com suas
responsabilidades bem delimitadas, são:

\begin{itemize}
	\item \texttt{app/core/extraction/twitter\_api.py}: implementa a classe \texttt{TwitterAPIExtractor}, que herda da classe abstrata \texttt{BaseExtractor} e é responsável pela construção e execução
	das requisições HTTP à API via biblioteca \texttt{requests}. Recebe como parâmetros o identificador do usuário (\texttt{x\_user\_id}), o intervalo temporal e, opcionalmente, o token de paginação;
	
	\item \texttt{app/core/extraction/base\_extractor.py}: define a classe abstrata \texttt{BaseExtractor} (ABC), que estabelece o contrato de interface para qualquer extrator de dados do 	\textit{pipeline}, garantindo extensibilidade para futuras fontes de dados além do X/Twitter;
	
	\item \texttt{app/shared/schemas.py}: módulo transversal a todas as etapas do \textit{pipeline}, contendo as \textit{dataclasses} de \textit{parsing} da resposta JSON da API — \texttt{Tweet}, \texttt{Tweets}, \texttt{NoteTweet}, \texttt{PublicMetrics}, \texttt{Entities}, \texttt{Hashtag} e \texttt{Meta} —, que estruturam os campos retornados em objetos Python tipados;
	
	\item \texttt{app/core/extraction/collection\_log.py}: define a classe \texttt{SearchTerm}, que representa um registro de coleta agendada com os campos \texttt{x\_user\_id}, 	\texttt{from\_date\_time} e \texttt{to\_date\_time}, e gerencia a inserção dos termos de busca na tabela \texttt{collection\_log};
	
	\item \texttt{app/shared/db/database.py}: módulo transversal que implementa a classe \texttt{DatabaseManager}, responsável pela interface de comunicação com o banco de dados PostgreSQL — incluindo as operações de consulta, inserção e atualização utilizadas por todos os estágios do \textit{pipeline};
	
	\item \texttt{app/core/extraction/\_\_main\_\_.py}: ponto de entrada do estágio de coleta (invocado via \texttt{python -m app.core.extraction}, que orquestra o ciclo completo: recupera
	registros com status \texttt{pending} em \texttt{collection\_log}, executa a paginação iterativa para cada um deles e persiste os resultados no banco de dados.
\end{itemize}

Os dados coletados foram persistidos em um banco de dados \textbf{PostgreSQL~15}, provisionado em contêiner \textbf{Docker} para garantir portabilidade e isolamento do ambiente. A infraestrutura é inicializada via \texttt{docker-compose.yml}, que sobe automaticamente o PostgreSQL (porta~5432) e o pgAdmin~4 (porta~5050), executando os scripts \texttt{infra/01-init-database.sh} — criação de usuário e permissões — e \texttt{infra/02-create-tables.sql} — criação das tabelas e índices. O esquema do banco é composto por três tabelas, cujas estruturas são descritas no \autoref{qua:schema_db}.

\begin{quadro}[!htb]
	\caption{\label{qua:schema_db}Estrutura das tabelas do banco de dados PostgreSQL}
	
	\begin{tabular}{|p{4.5cm}|p{4.5cm}|p{5.8cm}|}
		\hline
		\textbf{Tabela} & \textbf{Campo} & \textbf{Descrição} \\
		\hline
		\multirow{8}{*}{\texttt{tweets}} & \texttt{tweet\_id} & Identificador único da publicação na plataforma (\texttt{UNIQUE}) \\ & \texttt{username} & Identificador numérico do veículo autor \\
		& \texttt{note\_tweet} & Conteúdo textual utilizado na análise (\texttt{note\_tweet.text} ou \texttt{text}) \\ & \texttt{created\_at} & Data e hora de publicação (com fuso horário) \\ & \texttt{likes}
		& Número de curtidas \\ & \texttt{hashtags} & Lista de \textit{hashtags} (JSONB) \\ & \texttt{tweet} & Objeto JSON completo retornado pela API
		(JSONB) \\
		
		\hline
		\multirow{7}{*}{\texttt{tweets\_classification}} & \texttt{tweet\_id} & FK → \texttt{tweets.id} (\texttt{ON DELETE CASCADE}) \\ & \texttt{sentiment} & Polaridade atribuída: \textit{positivo}, \textit{negativo} ou \textit{neutro} \\ & \texttt{why\_sentiment} & Justificativa textual para a polaridade classificada \\ & \texttt{is\_finance\_news}   & Relevância financeira classificada: 1~=~financeiro, 0~=~não financeiro \\ & \texttt{why\_is\_finance\_news} & Justificativa textual para a classificação de relevância \\ & \texttt{classificator} & Identificador do classificador: \texttt{"Humano"}, \texttt{"FinBERT-PT-BR"}, \texttt{"BERTimbau"} ou outro modelo futuro \\
		
		\hline
		\multirow{6}{*}{\texttt{collection\_log}} & \texttt{search\_term} & Parâmetros da coleta agendada (JSONB) \\ & \texttt{tweets\_collected}    & Total de registros inseridos no 	ciclo \\
		& \texttt{start\_time} & Início da execução do ciclo de coleta \\ & \texttt{end\_time} & Término da execução do ciclo de coleta \\ 	& \texttt{status} & Estado do ciclo: \texttt{pending}, 		\texttt{completed} ou \texttt{partially\_completed} \\ & \texttt{error\_message} & Mensagem de erro, se houver \\
		\hline
	\end{tabular}
	\legend{Fonte: Elaborado pelo autor com base no schema \texttt{infra/02-create-tables.sql} do repositório do \textit{pipeline}.}
\end{quadro}

A tabela \texttt{collection\_log} cumpre função de controle de idempotência: cada combinação de veículo e intervalo temporal é registrada com status \texttt{pending} antes da execução, sendo
atualizada para \texttt{completed} ao final de um ciclo sem falhas de inserção, ou para \texttt{partially\_completed} caso algum registro não tenha sido persistido. Esse mecanismo permite a retomada segura de coletas interrompidas e a auditoria do histórico de execuções sem risco de duplicação de dados, uma vez que a coluna \texttt{tweet\_id} na tabela \texttt{tweets} é definida com restrição \texttt{UNIQUE}.
